(1) 可逆绝热理想气体满足 \(TV^{\gamma-1}=\mathrm{const}\)，于是 \(T_2=T_1r^{\gamma-1}\)，\(T_3=\Theta T_1r^{\gamma-1}\)，\(T_4=T_3/r^{\gamma-1}=\Theta T_1\)。在两段绝热过程，按气体对外做正功的符号，\(W_{12}=nC_V(T_1-T_2)<0\)、\(W_{34}=nC_V(T_3-T_4)>0\)。压缩过程需外界输入功，膨胀过程输出功。

(2) 定容段不做体积功，外界提供热 \(Q_{\rm in}=nC_V(T_3-T_2)=nC_V T_1r^{\gamma-1}(\Theta-1)\)。最后排热量的正数大小 \(Q_{\rm out}=nC_V(T_4-T_1)=nC_VT_1(\Theta-1)\)。全循环第一定律给 \(W_{\rm net}=Q_{\rm in}-Q_{\rm out}=nC_VT_1(\Theta-1)(r^{\gamma-1}-1)>0\)。这些表达式与两绝热段代数和完全一致。

(3) 热效率 \(\eta=W_{\rm net}/Q_{\rm in}=1-r^{1-\gamma}\)，不依赖本理想模型中的 \(\Theta\)。由于循环在定容吸、放热过程中气体温度连续变化，若仅使用固定高、低温热库接触，将不可逆并产生额外熵；即使定容过程设计为与连续多个热库可逆接触，也不能把其最高、最低温度代进单对热库 Carnot 循环公式来替代本 Otto 效率。

(4) 工质在绝热两段 \(\Delta S_{12}=\Delta S_{34}=0\)，定容吸热段 \(\Delta S_{23}=nC_V\ln(T_3/T_2)=nC_V\ln\Theta\)，排热段 \(\Delta S_{41}=-nC_V\ln\Theta\)，工质一周回到原态，\(\oint dS=0\)。若外部热库在每一瞬间与工质同温，则热库吸收的熵逐段相反，总熵产生可为零；实际仅两个恒温热库换热时应另按相应热库温度计入正的熵产生。
